\documentclass[11pt]{article}

\usepackage[utf8]{inputenc}
\usepackage[T1]{fontenc}
\usepackage[margin=0.4in, top=0.35in, bottom=0.4in, headheight=14pt]{geometry}
\usepackage{hyperref}
\usepackage{url}
\usepackage{booktabs}
\usepackage{amsfonts}
\usepackage{amsmath}
\usepackage{nicefrac}
\usepackage{xcolor}
\usepackage{titlesec}
\usepackage{enumitem}
\usepackage{fancyhdr}
\usepackage{graphicx}
\usepackage{subcaption}
\usepackage{caption}

\captionsetup{font=small, labelfont=bf, skip=2pt}
\captionsetup[subfigure]{font=footnotesize, labelfont=bf, skip=1pt}

\setlength{\parskip}{1pt}
\setlength{\parindent}{0pt}
\titlespacing*{\section}{0pt}{4pt}{1pt}
\titlespacing*{\paragraph}{0pt}{2pt}{4pt}

\titleformat{\section}{\normalfont\large\bfseries}{}{0em}{}
\titleformat{\paragraph}[runin]{\normalfont\normalsize\bfseries}{}{0em}{}[.\hspace{0.5em}]

\setlist[itemize]{noitemsep, topsep=1pt, leftmargin=1.5em}
\setlist[enumerate]{noitemsep, topsep=1pt, leftmargin=1.5em}

\renewcommand{\maketitle}{
  \begin{center}
    {\bf \LARGE {Predicting Showjumping Outcomes with Self-Supervised Video Representations}}\\
    \vspace{0.08cm}
    {\large CS 131 Project Milestone: Type-Conditioned Approach Quality from Unlabeled Footage}\\
    \vspace{0.08cm}
    {\bf Team Member:} Brooke Ballhaus \quad
    {\bf Email:} ballhaus@stanford.edu
  \end{center}
}

\begin{document}

\maketitle

\vspace{4pt}
\hrule
\vspace{8pt}

\section{Technical Progress}

I built the milestone-stage pipeline end-to-end on Colab Pro (A100): YouTube scraping, 2-second clip segmentation, YOLOv8 horse-presence filtering, hand-annotated fence boxes, geometric takeoff-distance extraction, and a self-supervised R(2+1)D-18 video encoder with InfoNCE contrastive and 6-way temporal-order objectives. The downstream classifier and regressor heads from the proposal are not yet trained --- those wait on the week-3 expansion of the labeled set. All results below use encoder embeddings, hand-drawn fence boxes, and the geometric $d$ estimator.

\paragraph{Data} I scraped 6 full Grand Prix highlight reels (FEI World Cup, Longines GCT, CSI 3*) at 1.50\,m, totaling roughly 12 hours of footage. A 30-second-stride window cutter produced $\mathbf{2{,}223}$ raw 2-second clips at 16\,fps, $320\times320$. A YOLOv8 horse-presence filter (keep clips with a horse visible in $\geq40\%$ of sampled frames) reduced this to $\mathbf{1{,}121}$ jumping clips and moved the rest --- crowd shots, course walks, slow-motion replays, interviews --- into a rejected pool for inspection. From these, I hand-annotated $\mathbf{39}$ fence bounding boxes (22 oxers, 17 verticals) using an inline matplotlib + ipywidgets annotator built for Colab; each box was drawn on a frame I scrubbed to manually rather than the clip midpoint, because broadcast cameras track the horse, not the fence.

\paragraph{Detection and geometric \texorpdfstring{$d$}{d}} Pretrained COCO YOLOv8 (the \texttt{horse} class) detects the horse every 2nd frame. The takeoff frame is found as the sharpest downward-to-upward reversal of the horse-box bottom edge in the second half of each clip --- a robust proxy for the hoof leaving the ground. The standard 3.5\,m top pole on the fence box gives meters-per-pixel as $\mathrm{mpp} = 3.5/w_{\mathrm{fence}}$. Takeoff distance $d$ is then the Euclidean ground distance between the takeoff-frame horse-box bottom-center (hoof line) and the fence base, scaled by mpp. Fence type (vertical vs.\ oxer) is currently a width-to-height heuristic with a pole-count override from the annotator, to be replaced in week 3 by a small CNN once $\geq 50$ typed labels are available.

\paragraph{Self-supervised pretraining} Encoder: torchvision R(2+1)D-18 trunk with a 2-layer MLP projector to 128-dim. Two pretext objectives are trained jointly: (a) InfoNCE on two photometric- and flip-augmented views of the same clip ($\tau=0.1$), with the rest of the batch as negatives, and (b) 6-way classification of the permutation applied to three disjoint 4-frame sub-clips, loss weight $\lambda_{\mathrm{order}}=0.3$. Optimizer: AdamW (lr $10^{-3}$, weight decay $10^{-4}$), cosine schedule, batch 16, 10 epochs on A100.

\section{Results and Visualizations}

\begin{figure}[!htb]
  \centering
  \begin{subfigure}{0.32\textwidth}
    \includegraphics[width=\linewidth]{figures/training_curve.png}
    \caption{SSL pretraining losses across 10 epochs. Both InfoNCE and temporal-order CE drop monotonically; order drops below the $\log 6 \approx 1.79$ chance floor.}
    \label{fig:loss}
  \end{subfigure}\hfill
  \begin{subfigure}{0.32\textwidth}
    \includegraphics[width=\linewidth]{figures/d_distribution.png}
    \caption{Geometric $d$ across the 39 annotated clips, split by fence type. Most values fall in the $1$--$3$\,m physically reasonable range; the $>3.5$\,m tail comes from oblique camera angles.}
    \label{fig:d}
  \end{subfigure}\hfill
  \begin{subfigure}{0.32\textwidth}
    \includegraphics[width=\linewidth]{figures/tsne_type.png}
    \caption{t-SNE of encoder embeddings colored by hand-annotated fence type. Colored points show no clean vertical/oxer separation; most points are grey (un-annotated).}
    \label{fig:tsne_type}
  \end{subfigure}
\end{figure}

\begin{figure}[!htb]
  \centering
  \begin{subfigure}{0.48\textwidth}
    \includegraphics[width=\linewidth]{figures/tsne_video.png}
    \caption{Same t-SNE, recolored by source broadcast. Clear per-video clusters --- the dominant axis of variation in the embedding space is venue (arena, lighting, footing, crowd), not approach quality. This identifies the main bottleneck for week 3.}
    \label{fig:tsne_video}
  \end{subfigure}\hfill
  \begin{subfigure}{0.48\textwidth}
    \includegraphics[width=\linewidth]{figures/det_grid.png}
    \caption{Four sample clips with horse box (green), hand-annotated fence box (orange), and the computed $d$ value overlaid at the detected takeoff frame. The bottom-right panel illustrates the oblique-angle failure mode.}
    \label{fig:det}
  \end{subfigure}
\end{figure}

\paragraph{What the figures say} Figure~\ref{fig:loss} confirms both SSL objectives optimize cleanly, so the encoder is genuinely picking up motion direction rather than appearance. Figure~\ref{fig:d} shows the geometric estimator produces values in the physically reasonable $1$--$3$\,m range. Figure~\ref{fig:tsne_type} shows the embedding does \emph{not} cleanly separate verticals from oxers --- but Figure~\ref{fig:tsne_video} explains why: the dominant variation is which broadcast a clip came from, not approach quality. This is the key empirical finding of the milestone. Figure~\ref{fig:det} shows representative geometric successes and the main failure mode ($d$ inflation under oblique angles).

\section{Deviations from the Proposal}

\paragraph{SSL collapsed from scratch} A 10-epoch run on 1{,}121 clips drove both losses down monotonically but yielded representational collapse (nearest-neighbor cosine similarity $\approx 0.97$ across unrelated clips). I executed the proposal's planned Kinetics-400 fallback --- the pretrained R(2+1)D-18 encoder used directly --- and the figures above all use these embeddings.

\paragraph{Venue confound} All 39 hand-annotated fences are from one video, picked for camera quality during a single annotation session. Combined with the venue-clustered embedding space, this blocks any approach-quality claim until the labeled set spans venues.

\paragraph{Fence detection} The proposal said YOLOv8 would label fence type from box shape and pole count; pretrained COCO YOLO has no fence class, so I hand-drew boxes for the milestone. Week 3 fine-tunes YOLO on these 39 boxes (plus new ones) to recover the original automated path.

\paragraph{Stretch goal deferred} Future-frame embedding prediction is not useful before the venue confound is fixed.

\section{Updated Timeline and Plan}

\begin{enumerate}
  \item \textbf{Week 3 (Cross-venue labels \& SSL fix).} Annotate $\sim 60$ more fence boxes spanning all 6 videos (target: 100 total). Fine-tune YOLOv8 on those boxes for automatic detection. Re-run SSL pretraining from Kinetics initialization with a venue-balanced sampler; add a domain-adversarial head to actively suppress venue identity. Train the proposal's three baselines (from-scratch supervised, ImageNet 2D-CNN, type-only logistic regression).
  \item \textbf{Week 4 (Analysis \& writeup).} Sample-efficiency curves at $n \in \{25, 50, 100, 200\}$, ablations on each pretext task and on type conditioning, $d$-MAE reported per fence type, t-SNE colored by outcome / $d$ / type, final report and presentation. Submit by the course deadline.
\end{enumerate}

\paragraph{Risks watched} (i) Cross-venue annotation is the rate-limiter; budgeted 90\,min/day. (ii) Domain-adversarial training can be unstable on small datasets --- I will fall back to a simpler arena-color jitter augmentation if it diverges. (iii) Side-view camera filtering will reduce the labeled-clip pool further; the clip-stride is tunable to compensate.

\end{document}
